\chapter{Úvod}
\label{chap:introduction}

Distribuované verzovací systémy, zejména systém Git, umožňují vývojářům asynchronně spolupracovat na sdílené codebase. Kolem nich vyrostly platformy jako GitHub, které dnes neslouží pouze jako úložiště zdrojových kódů, ale integrují i správu revizí kódu, sledování chyb a \CICD procesy.

Nativní analytické nástroje platformy GitHub dostačují pro malé a středně velké projekty, ale u velkých open-source repozitářů narážejí na své limity. Příkladem je repozitář překladače jazyka Rust (\texttt{rust-lang/rust}), kde spolupracují stovky vývojářů organizovaných do týmových hierarchií, otevírají se desítky \PR denně a stejně tak uzavírají.  Projekt využívá vlastní stavové automaty řízené štítky a striktní pravidla schvalování vynucovaná automatizovanými boty. Zjistit například, jak dlouho průměrně čeká kód na revizi od konkrétního týmu, nebo identifikovat stagnující \PR, je pomocí standardních nástrojů obtížné.

Cílem této diplomové práce je proto návrh a implementace analytického nástroje, který tyto informační mezery překlenuje. Navržený systém stahuje, normalizuje a analyzuje data ze tří zdrojů: z verzovacího systému Git, z \API platformy GitHub a z externích organizačních struktur (Team \API). Nástroj je implementován v jazyce Rust s ohledem na výkon při zpracování velkých objemů dat.

\section*{Cíle práce}
Práce nejprve teoreticky analyzuje mechanismy pro získávání dat z platformy GitHub (\REST \API, GraphQL \API a Webhooky) a popisuje specifika procesního řízení ve velkých open-source repozitářích. Součástí je také zhodnocení nativních možností monitoringu platformy GitHub i specializovaných nástrojů třetích stran a identifikace jejich nedostatků v kontextu masivních projektů. Na základě těchto zjištění je navržena modulární softwarová architektura a databázový model schopný ukládat historická i inkrementální data. Výsledné řešení je vyhodnoceno na reálných datech z repozitáře \texttt{rust-lang/rust} prostřednictvím sady analytických dotazů zaměřených na hlubší porozumnění procesu revize kódu.

\section*{Struktura práce}
Teoretická část práce popisuje principy verzovacích systémů, \API platformy GitHub a obsahuje analýzu existujících monitorovacích nástrojů. Navazuje kapitola o návrhu systému s definicí požadavků, návrhem architektury a strukturou databáze. Implementační kapitola se věnuje technickým specifikům vývoje analytického nástroje. Práce je zakončena kapitolou o vyhodnocení, která prezentuje výsledky měření a analytická data ze zkoumaného repozitáře.